\subsection{Testing the smooth proper noncommutative Weil framework}
There is a precise finite-field noncommutative Weil conjecture in a
different framework: for a smooth proper dg category, cyclotomic
Frobenius acts on the two finite-dimensional groups $TP_0[1/p]$
and $TP_1[1/p]$, with expected conjugate moduli $1$ and $\sqrt q$.
Tabuada \cite[\S1, pp.~2--3; \S5.1, pp.~13--14; \S6.1,
pp.~15--16]{TabuadaWeil} specifies these objects and proves the
comparison with smooth proper schemes in Theorem~1.5.
This cyclotomic Frobenius is not defined by raising matrices to
their $p$th powers. Its periodicity variable is also not the GCT
quantum parameter. The following test concerns the actual algebra
already constructed, without identifying the two conjectural programs.

\begin{finding}[\fid{F39}: the actual collision fiber fails homological smoothness]\label{f:ncweil}
In characteristic $p>3$, the abstract specialization of the
rank-three algebra in \fid{F19},
\[
 A_1=k\times k\times M_2(B),\qquad B=k[\epsilon]/(\epsilon^2),
\]
is proper but not smooth as a dg algebra concentrated in degree zero.
The same holds for its category of perfect complexes. Therefore this
particular category is not an input to the cited smooth proper
noncommutative Weil conjecture.

In contrast, at a split finite-field parameter point where the two
center roots are distinct and the stated denominators are units,
the algebra is $k^2\times M_2(k)^2$. Its perfect category is Morita
equivalent to four rational points. Its normalized even Frobenius
spectrum consists of four ones and its odd spectrum is empty.
Replacing $A_1$ by its semisimple quotient instead gives three
points and discards the nilpotent collision data.
\end{finding}
\begin{proof}
Properness of the one-object dg algebra means that its underlying
complex over $k$ is perfect; this holds because $A_1$ is finite
dimensional. Homological smoothness would make its diagonal bimodule
perfect over $A_1\otimes_k A_1^{\rm op}$.
We test the dual-number block directly.
Put
\[
 S=B\otimes_k B=k[x,y]/(x^2,y^2),\qquad
 u=x-y,\quad w=x+y.
\]
There is a free $S$-resolution of the diagonal $B$,
\[
 \cdots\xrightarrow{u}S\xrightarrow{w}S
       \xrightarrow{u}S\longrightarrow B\longrightarrow0.
\]
Here the rightmost positive differential is multiplication by $u$.
To check exactness in full, use the basis $1,x,y,xy$ of $S$.
The kernel of the augmentation $x,y\mapsto\epsilon$ is
$\langle x-y,xy\rangle=\operatorname{im}(u)$.
Multiplication by $u$ sends
$a+bx+cy+dxy$ to $a(x-y)+(c-b)xy$, so its kernel is
$\langle x+y,xy\rangle=\operatorname{im}(w)$.
Multiplication by $w$ has kernel
$\langle x-y,xy\rangle=\operatorname{im}(u)$.
Also $uw=0$. These identities prove exactness at every stage.

After tensoring over $S$ with the diagonal $B$, the differentials
alternate between zero and multiplication by $2\epsilon$. Hence
\[
 HH_0(B)=B,\qquad
 HH_{2j+1}(B)=B/(\epsilon)\ (j\ge0),\qquad
 HH_{2j}(B)=k\epsilon\ (j\ge1).
\]
In particular $\operatorname{Tor}^{S}_i(B,B)$ is nonzero for
every $i>0$. A perfect $S$-module would have bounded derived tensor
product with $B$, so $B$ is not homologically smooth.
If $A_1$ were smooth, restricting its diagonal along the central
matrix-block idempotent would make the diagonal of $M_2(B)$
perfect. The matrix Morita equivalence, obtained by the full matrix
unit $E_{11}$ on each side, is exact and preserves finite projective
resolutions; it would then make $B$ perfect over $S$, a contradiction.
This also proves the assertion for the Morita-invariant perfect dg
category, using the smoothness definition in the cited source.

At a split good parameter the distinct-root Chinese remainder
decomposition of \fid{F19} gives $k^2\times M_2(k)^2$.
Applying ordinary matrix Morita equivalence to its two matrix factors
identifies its perfect category with that of $\coprod_1^4\Spec k$.
The latter is smooth proper. The crystalline comparison
\cite[\S7, equations~(7.1)--(7.4), p.~17]{TabuadaWeil}
has only four degree-zero classes, each with Frobenius one; this
proves the stated normalized spectrum. The radical of $A_1$ is
the square-zero matrix ideal $M_2(k\epsilon)$, so its quotient is
$k^2\times M_2(k)$ and the same calculation gives three points.
\end{proof}

The dual-number resolution is the familiar elementary computation;
its use here is a test of the actual nonstandard collision fiber,
not a new claim about dual numbers or priority over noncommutative
motivic theory. A smooth categorical replacement or a relative theory
is not excluded. Such a replacement would need a specified comparison
retaining the coherent specialization and the designated based
operations. The pure spectrum at the good fibers does not supply
that comparison, and semisimplification expressly loses data measured
by \fid{F25}. The smoothness failure also does not refute any
conjecture whose domain already requires smoothness.

% Separate paragraph for Part III.
\paragraph{\fid{F39} in Grothendieck's language.}
The singular piece is the actual affine scheme
$\Spec k[\epsilon]/\epsilon^2$, and matrix Morita equivalence
expresses its coherent category without changing the defect.
Smoothness here asks for a perfect diagonal bimodule, a perfect-complex
condition distinct from finite-dimensionality or ordinary $\ell$-adic
purity. The displayed resolution proves its failure directly. EGA~I's
affine scheme language and the perfect-complex framework of SGA~6
provide the classical dictionary; the smooth proper dg and
noncommutative-motive extension is explicitly the modern cited theory.
The audited cote~16, scans~48--49, concerns conditional semisimplicity
and Weil statements in its own motivic framework; it is not a source
for this dg construction. A new realization and basis comparison is
required before its Frobenius can address the GCT coefficients.
